\providecommand{\vTwoNumberMacros}{paper1_v2_high_severity_draft/tables/multiseverity_numbers}
\input{\vTwoNumberMacros}

\subsection{IR and SR across a Wider Severity Range}
\label{sec:exp-multiseverity-extension}

The original narrow Gaussian experiment evaluates a moderate visual shift,
$\sigma_{\mathrm{eval}}=\vTwoEvalZeroEight$, where every checkpoint passing the
IR condition also passes SR. Does the SR condition become more active when the
evaluated shift becomes more severe? We address this question in a separate
experiment rather than extending the original threshold-selection sweep. For each of
\vTwoTaskCount{} tasks and \vTwoRunCount{} training runs, we train checkpoints with
$\sigma_{\max}\in\vTwoTrainSigmaSet$ and evaluate them at
$\sigma_{\mathrm{eval}}\in\vTwoEvalSigmaSet$. The eight-step rollout
protocol, state-pair catalog, and thresholds
$(t_{\mathrm{IR}},t_{\mathrm{SR}})=(\vTwoIRThreshold,\vTwoSRThreshold)$ are
unchanged.

\begin{figure}[t]
\centering
\includegraphics[width=0.82\linewidth]{fig_sr_guard_activation_by_severity.pdf}
\caption{\textbf{IR and SR become restrictive at different evaluation severities.}
Both curves use the same denominator of \vTwoCheckpointCount{}
task--training-run--training-strength rows
($\vTwoTaskCount\times\vTwoRunCount\times\vTwoTrainLevelCount$) at each
severity. The blue curve counts rows passing IR; the orange curve counts rows
passing both conditions, and the shaded gap is the additional SR veto. The
IR-pass counts are \vTwoIRPassAtZeroEight{}, \vTwoIRPassAtZeroSixteen{},
\vTwoIRPassAtZeroThirtyTwo{}, and \vTwoIRPassAtZeroSixtyFour{} as
$\sigma_{\mathrm{eval}}$ increases; SR vetoes \vTwoVetoAtZeroEight{},
\vTwoVetoAtZeroSixteen{}, \vTwoVetoAtZeroThirtyTwo{}, and
\vTwoVetoAtZeroSixtyFour{} of those rows, leaving
\vTwoJointPassAtZeroEight{}, \vTwoJointPassAtZeroSixteen{},
\vTwoJointPassAtZeroThirtyTwo{}, and \vTwoJointPassAtZeroSixtyFour{} joint
passes. Counts are descriptive: the same checkpoints recur across severity
slices, so the points are not independent; the task breakdown is in
\Cref{sec:appendix-multiseverity}. This decomposition shows when each condition
is active; it does not measure checkpoint-selection accuracy.}
\label{fig:sr-guard-activation}
\end{figure}

The low-severity failures in \Cref{fig:sr-guard-activation} are the same
\vTwoVetoAtZeroEight{} TwoRoom checkpoints: \vTwoLowVetoAtOne{} trained with
$\sigma_{\max}=\vTwoTrainOne$ and \vTwoLowVetoAtOneFive{} with
$\sigma_{\max}=\vTwoTrainOneFive$. These training levels are absent from the
original narrow grid, so the low-severity vetoes do not contradict that result.
At higher evaluation severities, failures extend to PushT and one Reacher
checkpoint; the full breakdown is reported in
\Cref{sec:appendix-multiseverity}.

We next ask whether the continuous diagnostics vary with same-checkpoint
planning-performance retention, not merely whether they cross a threshold.
For each checkpoint, we define
retained planning performance as success under noise minus success on clean
observations for that same checkpoint. Within each task and training run, we
remove the additive training-strength and evaluation-severity trends from both
the diagnostics and retained performance, then compute rank associations over
the remaining $\vTwoTrainLevelCount\times\vTwoEvalLevelCount$ interaction.
Training runs are averaged within each
task before the \vTwoTaskCount{} tasks receive equal weight.

The direct equal-task mean Spearman correlations are similar in magnitude:
$\vTwoIRDirect$ for lower relative IR, $\vTwoSRDirect$ for higher SR, and
$\vTwoExposureDirect$ for the simple exposure ratio
$-\sigma_{\mathrm{eval}}/\sigma_{\max}$. We do not interpret the small
differences among these direct summaries as evidence of diagnostic superiority.
After rank-residualizing the diagnostics and retained performance against the
exposure ratio, the correlations remain positive at $\vTwoIRAdjusted$ for
lower relative IR and $\vTwoSRAdjusted$ for higher SR. The corresponding task
means are positive for all four tasks; ranges and a retention-ratio sensitivity
analysis appear in \Cref{sec:appendix-multiseverity}.

These results clarify SR's regime-dependent role. In the original
moderate-severity sweep, IR tracks recovery while SR checks whether the tested
state-coordinate distinctions remain outside the contracted radius.
In the broader experiment, the SR-veto fraction among IR-passing rows is higher
in the harder slices, and the continuous SR values remain associated with
same-checkpoint retention beyond a simple exposure ratio. This evidence
does not show that the two-condition screen selects better checkpoints. It also
does not replace nominal model evaluation: a checkpoint can preserve its
relative performance under noise while performing poorly on both clean and
noisy inputs.
\FloatBarrier
